\documentclass[a4paper]{article}
\usepackage[utf8]{inputenc}

\usepackage{graphicx}
\usepackage[bf]{subfigure}
\usepackage[bf]{caption}
% Section titles in sans-serif:
\usepackage{sectsty}
\usepackage{abstract}
\allsectionsfont{\sffamily}

%\renewcommand{\abstractname}{\textsf{Abstract} xxxx}
\renewcommand{\abstractnamefont}{\bfseries\textsf}

\usepackage[usenames,dvipsnames]{color}
\definecolor{seagreen}{RGB}{46,139,87}
\definecolor{codebg}{RGB}{255,255,150}
\usepackage{listings}
\usepackage[colorlinks]{hyperref}
\usepackage{url}	

\hypersetup{%
  citecolor=blue,
  linkcolor=blue,
  urlcolor=blue,
}%



\lstset{
  extendedchars=\true,
  inputencoding=utf8,
  language=Matlab,
  %showstringspaces=false,
  formfeed=\newpage,
  tabsize=4,
  %commentstyle=\itshape,
  basicstyle=\ttfamily\scriptsize,
  %basicstyle={\small\fontfamily{fvm}\fontseries{m}\selectfont},
  commentstyle=\color{Apricot}\bfseries,
  %commentstyle=\color{red}\itshape,
  stringstyle=\color{red},
  identifierstyle=\color{PineGreen},
  showstringspaces=false,
  keywordstyle=\color{blue}\bfseries,
  moredelim=[il][\large\textbf]{\#\# },
  morekeywords={models,range},
  numbers=left,
  numbersep=2pt,
  numberstyle=\tiny,%\color{blue}\bfseries,
  backgroundcolor=\color{codebg},
  literate=%
  {ã}{{\~a}}1
  {â}{{\^a}}1
  {õ}{{\~o}}1
  {á}{{\'a}}1
  {ú}{{\'u}}1
  {í}{{\'i}}1
  {é}{{\'e}}1
  {Ç}{{\c{C}}}1
  {Õ}{{\~O}}1
  {Ê}{{\^E}}1
  {ó}{{\'o}}1
  {à}{{\`a}}1
  {Â}{{\^A}}1
  {ô}{{\^o}}1
  {ê}{{\^e}}1
  {ç}{{\c{c}}}1
}

\newcommand{\code}[2]{
 \vspace{1em}
 \subsubsection*{#1}
 \lstinputlisting{#2}
}

%%%%%%%%%%%%%%%%%%%%%%%%%%%%%%%%%%%%%%%%
% You have two versions of the macro
% \draftnote{My note}. The first version puts notes (e.g. My note in the example)
% into the margin of your document. The second formats the note as nothing. You
% 'comment out' the version of the macro you don't want (using a % at the
% beginning of the line).
\newcommand{\draftnote}[1]{\marginpar{\tiny\raggedright\textsf{\hspace{0pt}#1}}}
%\newcommand{\draftnote}[1]{}

% This one is just for the comments for in-line text.
\newcommand{\indraftnote}[1]{\textcolor[HTML]{114406}{\texttt{\footnotesize[#1]}}}
%\newcommand{\indraftnote}[1]{}
\newcommand{\todo}[1]{\indraftnote{todo: #1}}
\newcommand{\ie}{{\it i.e.}}
\newcommand{\etc}{{\it etc}}
\newcommand{\eg}{{\it e.g.}}
\newcommand{\wrt}{{\it w.r.t. }}
\newcommand{\etal}{{\it et.\ al.\ }}
\newcommand{\etalf}{{\it et.\ al.}}


\begin{document}

\title{\textsf{Computer Graphics Lab #2\\ Animation and 3D}
\marginpar{\vspace{-1.2cm}\includegraphics[height=1.4cm] {figs/logo_uerj_cor_small.png}}} 

\author{Prof.\ Ricardo Fabbri, Ph.D.\\[1em]
\small{
Polytechnic Institute (IPRJ) at the Rio de Janeiro State University\\
\url{http://wiki.nosdigitais.teia.org.br/CG}}
}
 

\date{\today}
\maketitle
\begin{abstract}
Now that you did Lab\#1 which consisted of rotating a camera around an
arbitrary axis, its time to move on to further concepts and consolidate your
learning. This should help you with the bigger goal of the final project
of generating a complex animation.
\end{abstract}
\vspace{2em}

Any material need for the lab can be found from the course
website. \textbf{All extra work will be considered for bonus grade points.}

\section{Construct the IPRJ Logo in 2D}
Your goal is to generate the IPRJ logo of Figure~\ref{fig:iprj:logo} in Blender,
in multiple ways. Steps:
\begin{enumerate}
\item Construct the IPRJ logo in 2D using the standard Blender API, without any
data file. The standard blender API is used to help you construct models that
can be further manipulated by other Blender scripts or GUI. Use 
spline curves and surfaces for the sine wave and curved letters, possibly for everything.
\item Construct the IPRJ logo in 2D using the Blender OpenGL API, without any
data file. The OpenGL API is best for making blender tools that display
interactive graphics
\emph{prior} to a full render. Use 
spline curves and surfaces for the sine wave and curved letters, possibly for everything.
\end{enumerate}

\section{Animate the IPRJ Logo}
\begin{enumerate}
\item Create an animation of the IPRJ logo. Make the sine wave oscilate in 2D or 
3D (off the plane). 
\item Make the ball bounce. 
\item Devise your own effects. Be creative! If done with 
stunning, professional high quality, this can pass as a final project.
\end{enumerate}

\bibliographystyle{IEEEtran}
\bibliography{personal}

\end{document}
